\documentclass[11pt]{article}

\usepackage[margin=0.82in]{geometry}
\usepackage{amsmath}
\usepackage{booktabs}
\usepackage{array}
\usepackage{tabularx}
\usepackage{microtype}
\usepackage{needspace}
\usepackage{xcolor}
\usepackage{xurl}
\usepackage[hidelinks]{hyperref}
\usepackage{placeins}
\usepackage{pgfplots}
\usepgfplotslibrary{groupplots}
\pgfplotsset{compat=1.18, /pgf/number format/1000 sep={}}

\setlength{\parindent}{0pt}
\setlength{\parskip}{0.55em}
\definecolor{stressred}{RGB}{184,48,62}
\definecolor{debtblue}{RGB}{48,105,152}
\definecolor{policygold}{RGB}{185,124,24}
\definecolor{quietgray}{RGB}{95,95,95}
\newcolumntype{Y}{>{\raggedright\arraybackslash}X}
\hypersetup{
  pdftitle={What Eventually Forces Federal Deficits to Change?},
  pdfauthor={U.S. Fiscal Dynamics Simulator}
}

\title{What Eventually Forces Federal Deficits to Change?\\
\large A Review of Fiscal Sustainability and the Confidence-Shock Simulation}
\author{U.S. Fiscal Dynamics Simulator}
\date{September 5, 2026}

\begin{document}
\maketitle

\begin{quote}
\small
\textbf{Companion review.} This September 5 review is preserved separately from
the \href{confidence_shock_household_scenario.pdf}{restored mechanism-led scenario}.
It evaluates that scenario's economic assumptions; it does not replace the worked
example. A \href{continuing_to_borrow.pdf}{self-contained overview} explains three
findings that survive the review. The
\href{confidence_shock_household_scenario_2026-08-29.pdf}{untouched August 29 edition}
is retained for comparison. Document links and this note were updated September 7, 2026;
the review's analysis and numerical results are unchanged.
\end{quote}

\begin{abstract}
Persistent federal deficits do not imply an inevitable bankruptcy date. A government can
roll over debt indefinitely when its primary balance, borrowing costs, economic growth,
and demand for its liabilities are mutually consistent. The U.S. problem in the pinned
February 2026 baseline is that debt keeps growing relative to income, leaving an increasing
claim on future budgets and less room to absorb adverse developments. If that imbalance
persists beyond what investors will hold at returns compatible with stable purchasing
power, some combination of fiscal promises, real creditor returns, or the economic
environment must change. Accounting alone cannot select the adjustment.

This review retains the Treasury cohort engine but revises the interpretation of the
500-basis-point confidence experiment. Its shock is imposed, its crisis date comes from an
unestimated gross-issuance ceiling, and its inflation spiral depends strongly on assumed
macroeconomic transmission. Changing quarterly inflation persistence from 0.98 to 0.90,
with the same shock and bondholder spending, reduces 2056 inflation from 8.04 to 2.89 percent
while debt reaches 260 percent of GDP. The fiscal deterioration is much more robust than
the proposed inflationary endgame.

New no-shock calculations find debt/GDP of 172.77 percent in 2056Q3 and a final-year increase
of 3.13 percentage points. Under unchanged reference growth and rates, a permanent
primary-balance improvement of 2.68 percent of GDP beginning in 2026Q4 satisfies a stated
2056 debt ceiling and final-year nonincrease condition; a four-year phase-in requires
2.70 percent. These are conditional fiscal calculations, not forecasts or optimal policies.
The answer is a constraint on the eventual adjustment, not a prediction of a particular
crisis, inflation rate, or date.
\end{abstract}

\section{The answer to the original question}

The United States need not pay off its federal debt, and it need not eliminate every annual
deficit. It does need a way to keep the claims it issues compatible with the economy's
capacity to support them and the public's willingness to hold them. Borrowing can postpone
an adjustment and spread costs across time. It does not by itself create the future goods,
services, or taxable income against which the claims are issued.

The distinction between the \emph{total deficit} and the \emph{primary deficit} matters.
The total includes interest; the primary deficit excludes it. Even a government running
primary surpluses can maintain a growing nominal debt stock and recurring total deficits
while its debt/GDP ratio falls. Conversely, a primary deficit can be sustainable if the
interest rate is sufficiently below the economy's growth rate.

What makes the U.S. baseline troubling is a continuing increase in debt relative to income,
under assumptions that already allow ongoing refinancing. The practical consequences can
start well before any market crisis: larger interest commitments, pressure on future taxes
and benefits, exposure to refinancing rates, and less flexibility during recessions. Higher
public borrowing can also displace private investment; the size depends on saving and
international capital flows. These are economic costs without requiring a failed auction
\cite{bernankeSustainability}.

\textbf{If the fiscal gap remains too large to be absorbed by growth and voluntary demand
for public liabilities, the existing promises cannot all be delivered on their original
real terms.} Taxes or spending must adjust, creditors must receive lower real returns,
payments must change, or favorable growth and financing conditions must make the gap
manageable. Inflation, financial repression, default, and fiscal reform are different
ways of allocating that adjustment. The simulation does not establish which will prevail.

This is a substantive answer to ``why is it unsustainable?'' It is also a limit on what
``what finally happens'' can mean. If every fiscal, monetary, regulatory, payment, and
economic assumption is fixed in a configuration that cannot be financed indefinitely,
there is no internally consistent indefinite continuation to forecast. A model must specify
which assumption gives way. An arbitrary numerical ceiling cannot supply that missing
political and economic behavior.

\section{What does ``no policy changes'' mean?}

The main counterfactual here is \emph{continued scheduled fiscal commitments and the revenue
path in the pinned CBO baseline}, with additional borrowing allowed. It is not literally
Congress doing nothing forever. CBO's February 2026 report assumes scheduled entitlement
payments continue after trust-fund exhaustion under its baseline conventions. It also
explains that authority to pay the full amounts would then be absent without legislation.
In that vintage, OASI exhaustion occurs in 2032; an alternative limited to available income
has materially lower benefits \cite{cbo2026}.

Literal legislative inaction therefore has a concrete possible consequence before any
modeled bond-market event: the government cannot deliver all scheduled benefits under the
existing financing authority. This paper does not simulate that payable-benefits alternative
or assert a precise payment procedure after exhaustion. It uses continued scheduled
commitments to examine the financing problem behind the promises. It likewise abstracts
from a binding statutory debt limit. Legal authority to issue debt and economic willingness
to hold debt are separate constraints.

The fiscal inputs end in 2056. An assertion about unchanged policy \emph{forever} would also
need a post-2056 path for demographics, health costs, taxes, spending, productivity, and
interest rates. The illustrative infinite-horizon arithmetic below explicitly states such
assumptions; it is not an extension of CBO's forecast.

\section{The debt arithmetic, including a sustainable-deficit counterexample}

Let $B_t$ be end-of-year nominal debt, $Y_t$ nominal GDP, $i_t$ the effective nominal
interest rate on opening debt, $D_t$ the primary deficit, and $A_t$ other financing
adjustments. Define $b_t=B_t/Y_t$, $d_t=D_t/Y_t$, $a_t=A_t/Y_t$, and nominal GDP growth
$n_t=Y_t/Y_{t-1}-1$. Then
\begin{align}
B_t &= (1+i_t)B_{t-1}+D_t+A_t,\\
b_t-b_{t-1} &= \frac{i_t-n_t}{1+n_t}b_{t-1}+d_t+a_t.
\label{eq:debt}
\end{align}
This is an accounting identity, with effective interest defined consistently with the debt
measure. Gross refinancing is absent: repaying principal with replacement debt does not
increase the stock. In the cohort engine, interest includes TIPS principal compensation;
other financing includes any explicit buyback premium or discount.

The primary deficit consistent with a locally unchanged debt ratio is
\[
d_t^{\rm stable}=\frac{n_t-i_t}{1+n_t}b_{t-1}-a_t.
\]
At debt of 100 percent of GDP and nominal growth of 4 percent, a 3-percent effective
interest rate permits a primary deficit of about 0.96 percent of GDP. A 5-percent rate
instead requires a primary surplus of about 0.96 percent. Nominal rates must be compared
with nominal growth; real rates with real growth using consistent compounding.

For constant $i,n,d$, with $a=0$, the recursion is $b_t=q b_{t-1}+d$, where
$q=(1+i)/(1+n)$. If $i<n$, it converges to
\[
b^*=\frac{d(1+n)}{n-i}.
\]
For $i=n$ and $d>0$, debt/GDP rises linearly; for $i>n$ and $d>0$, it grows without bound.
These are statements about the specified recursion, not about when a market must panic.

\begin{center}
\small
\begin{tabular}{lrrrr}
\toprule
Illustration & Primary deficit & Interest & Year-30 debt/GDP & Long-run limit\\
\midrule
Small deficit, favorable rates & 1.0\% & 3.0\% & 101.01\% & 104\%\\
Larger deficit, favorable rates & 2.5\% & 3.0\% & 140.26\% & 260\%\\
Interest equals growth & 2.5\% & 4.0\% & 175.00\% & unbounded\\
Interest exceeds growth & 2.5\% & 5.0\% & 219.71\% & unbounded\\
\bottomrule
\end{tabular}
\end{center}
All rows start at 100-percent debt/GDP, hold nominal growth at 4 percent, and have no
financing adjustments. The first two are counterexamples to the claim that every permanent
primary deficit must end in crisis. Their fixed low rates, even as debt rises, are strong
assumptions. Blanchard's analysis explains why low safe interest rates can permit rollover
without later tax increases, while distinguishing fiscal feasibility from welfare costs
\cite{blanchard}. A favorable historical average rate--growth differential does not insure
against future shocks or guarantee the rate offered on additional debt.

\section{What happens in this model without a confidence shock?}

The reconstructed Treasury stock enters 2026Q4 at approximately \$32.2 trillion, or
99.22 percent of opening GDP. Bills, notes, bonds, TIPS, floating-rate notes, and a coarse
other-public-debt component retain their existing accounting treatment. Inherited fixed-rate
coupons do not jump when marginal rates change. The primary path, GDP, and issuance-rate
reference use the pinned February 2026 data and documented quarterly interpolation.

The new no-shock run reproduces the existing reference path: 2056Q3 debt/GDP is
172.77 percent and inflation is 2.06 percent. Using average fiscal-year GDP, the underlying
baseline validation reports 175.06-percent debt/GDP, compared with CBO's 175.08 percent.
The difference between 172.77 and 175.06 is the GDP denominator convention, not an economic
disagreement. The validation is documented in \path{docs/validation.md}.

To expose the imbalance, the new diagnostic decomposes the last four quarterly changes
in the no-shock debt ratio. Each contribution uses its own quarter's GDP, so the sum
reconciles exactly to the change from 2055Q3 to 2056Q3:
\begin{center}
\begin{tabular}{lr}
\toprule
Contribution & Percentage points of debt/GDP\\
\midrule
Primary deficits & $+2.20$\\
Modeled interest cost & $+6.88$\\
Nominal GDP growth & $-5.95$\\
Other financing & $0.00$\\
\midrule
Change in debt/GDP & $+3.13$\\
\bottomrule
\end{tabular}
\end{center}
This is the mechanism requiring attention before a panic is introduced. Growth absorbs
part of the interest and primary borrowing, but the combined flow still raises debt/GDP.
The figures are contributions to a ratio change, not fiscal-year flow shares calculated
with one annual GDP denominator. Modeled interest is the cohort financing cost, not CBO's
net-interest budget concept.

The 3.13-point increase does not imply a crisis in 2056 or prove divergence under every
possible later path. It shows that this baseline has not stabilized by its endpoint.
Nor is 3.13 automatically the permanent fiscal package needed today: an earlier adjustment
changes future borrowing and therefore future interest expense.

\subsection*{How much preventive adjustment does the cohort model require?}

The new solver exercise holds reference growth, inflation, and issuance rates fixed. It
reduces the primary deficit permanently by a constant share of contemporaneous GDP,
starting in 2026Q4. Success requires 2056Q3 debt/GDP no greater than its initial
99.22 percent and no increase over the final four quarters. The trend condition has zero
tolerance apart from numerical solver precision.

\begin{center}
\begin{tabular}{lrrr}
\toprule
Implementation & Primary improvement & 2056 debt/GDP & Final-year trend\\
\midrule
Immediate & 2.68\% of GDP & 86.72\% & approximately flat\\
Four-year phase-in & 2.70\% of GDP & 91.54\% & approximately flat\\
\bottomrule
\end{tabular}
\end{center}
The final-year condition binds before the debt ceiling does. Consequently, the terminal
ratios are below 99.22 percent. This is not a recommendation to target 87 or 92 percent,
and it is not a proof of stability after 2056. Changing the horizon or required terminal
trend changes the solution. The calculation also omits the package's effects on output,
taxable income, interest rates, and political durability.

\begin{figure}[htbp]
\centering
\begin{tikzpicture}
\begin{axis}[
  width=0.91\textwidth, height=0.43\textwidth,
  xmin=2026.5, xmax=2057, ymin=70, ymax=185,
  xtick={2027,2036,2046,2056}, ylabel={Debt held by the public, percent of GDP},
  grid=major, grid style={draw=gray!25},
  legend style={font=\small, draw=none}, legend pos=north west
]
\addplot[debtblue, very thick] table[col sep=comma, x=plot_year,
  y=baseline_debt_pct] {sustainability_review_plot_data.csv};
\addlegendentry{No-shock reference}
\addplot[policygold, very thick] table[col sep=comma, x=plot_year,
  y=immediate_adjustment_debt_pct] {sustainability_review_plot_data.csv};
\addlegendentry{Immediate primary improvement: 2.68\% of GDP}
\addplot[stressred, dashed, very thick] table[col sep=comma, x=plot_year,
  y=phased_adjustment_debt_pct] {sustainability_review_plot_data.csv};
\addlegendentry{Four-year phase-in: 2.70\% of GDP}
\end{axis}
\end{tikzpicture}
\caption{Preventive fiscal arithmetic without a confidence shock. All three paths use
the same reference growth, inflation, and marginal rates. The adjustment paths satisfy
a finite-horizon endpoint and final-year trend condition; they are not macroeconomic
forecasts of tax increases or spending reductions.}
\label{fig:no-shock}
\end{figure}
\FloatBarrier

\section{Where the economic constraint comes from}

\subsection*{A buyer can always be found only under additional conditions}

A Treasury auction exchanges assets; it does not consume a fixed national pile of savings.
Maturing bonds return principal to investors that can be reinvested. Deficits add public
liabilities and corresponding financial claims elsewhere in the economy. Banks, funds,
households, foreign investors, and the Federal Reserve can change their portfolios.
Consequently, gross issuance cannot be compared with an assumed stock of ``money left''
and interpreted as a national solvency test.

What matters is the desired holding of public liabilities at their offered prices and
returns, alongside the real resources and fiscal responses that support them. Treasury
securities provide collateral, liquidity, and safe-asset services. Those services can
support unusually low borrowing costs. They do not guarantee unchanged demand at every
possible debt level or under every monetary and fiscal regime.

For orientation, a deterministic one-period real-debt identity can be iterated as
\[
L_0=\sum_{t=1}^{T}\frac{S_t}{\prod_{j=1}^{t}(1+r_j)}
       +\frac{L_T}{\prod_{j=1}^{T}(1+r_j)},
\]
where $L$ is real debt and $S$ the primary surplus, with other financing omitted.
The terminal liability is explicit. Turning this into an infinite-horizon statement that
debt equals discounted future surpluses requires a terminal condition. That condition
cannot be silently imposed while simultaneously claiming to test the feasibility of
indefinite rollover. Risk, liquidity services, and monetary liabilities require additional
structure beyond this simple identity.

\subsection*{Higher yields can postpone the problem and also worsen it}

If investors require a higher return, the price of existing fixed-rate bonds falls and
newly issued securities promise higher coupons. The capital loss occurs promptly; the
Treasury cash-flow effect arrives as debt matures. If primary balances do not respond,
higher future interest must be financed by further borrowing. If higher financing costs
also depress investment and the tax base, the fiscal gap can widen on both sides.

Bhatt and coauthors estimate local effects of expected debt on longer-run neutral rates
and the ten-year term premium, roughly 1--2 and 2--3 basis points respectively per
percentage point of debt/GDP \cite{fedDebtRates}. Those estimates support a possible feedback,
not a crisis threshold or unlimited linear extrapolation. A debt supply change, a change
in inflation expectations, and a liquidity panic need not have the same effects on either
Treasury yields or private borrowing rates.

The underlying conflict emerges when the additional return required by investors makes
the unchanged fiscal promises less credible, instead of restoring a sustainable borrowing
path. Expectations about later taxes, benefits, inflation, and payment behavior can then
move asset prices today. There need not be a clean sequence in which a quantity limit is
first reached and policymakers only then begin to react.

\section{What the confidence-shock model establishes, and what it does not}

The previous version built a vivid history around a ten-year, 500-basis-point premium on
new Treasury issuance, followed by an assumed financing boundary and a fiscal settlement.
The Treasury accounting remains useful. The narrative overstated what that accounting and
the chosen macro rules identify. The invented news chronology has been removed from this
research report; conditional stress and recovery calculations remain available below.

\subsection*{The no-shock path is protected by construction}

Incremental interest demand is measured relative to a parallel no-shock ledger. The debt
premium depends on the gap from that ledger's debt/GDP ratio. Financing capacity is
normalized to its gross issuance, scaled to scenario GDP. With no initiating disturbance,
all three mechanisms reproduce the reference even while baseline debt rises.

The new review verifies this directly: turning on both the central debt--yield rule and
the central capacity rule without imposing a shock leaves the entire no-shock result
unchanged. Avoiding double counting of baseline effects is a reasonable modeling choice.
But it means this engine cannot discover that the baseline itself eventually loses fiscal
credibility. A model designed to answer that question needs an independently justified
fiscal backing or investor-demand mechanism.

\subsection*{Gross interest received is not a sufficient statistic for net spending}

The earlier central allocation sends 15.67 percent of additional cash interest into
domestic demand: 56.88-percent domestic private holdings times a 25-percent spending
fraction, plus 28.97-percent foreign holdings times a 5-percent fraction. The remaining
14.15-percent Federal Reserve allocation contributes zero directly.

These holdings combine projected CBO aggregates with a TIC foreign-holdings benchmark
\cite{tic}. Holdings shares are not interest-income shares when maturities and yields
differ across holders. Domestic holdings include intermediaries and other institutions;
they are not a household consumption survey. Taxes on interest, pension retention,
deposit pass-through, wealth losses on existing securities, borrowers' interest exposure,
and expectations of future fiscal adjustment affect the aggregate demand response.

Auclert's redistribution framework explicitly considers winners and losers with different
spending propensities; its empirical sufficient statistics suggest amplification of monetary
policy transmission. It does not establish that rate increases become inflationary at
this paper's selected spending share \cite{auclert}. The BIS discusses an offsetting
government-interest-income channel, but also bond valuation losses and fiscal-risk
repricing, with different effects on output and inflation \cite{bisDebt2026}.

The interest-income channel is economically possible. Treating its assumed magnitude as
evidence that the Fed necessarily loses control is not justified. The reference-relative
GDP-share comparison avoids a nominal-price-level arithmetic error; it does not identify
a causal aggregate consumption response.

\subsection*{The shock's contractionary transmission is incomplete}

In the original run, the five-point Treasury premium raises federal financing costs for
forty quarters. It does not enter private-demand restraint directly. Only the Fed's
real policy-rate adjustment and a separate two-point private-credit spread lasting eight
quarters do so. Thus the experiment gives the Treasury premium a decade of fiscal cost
without a corresponding decade of private borrowing-cost or asset-valuation effects.
The correct pass-through depends on the shock; zero cannot be treated as innocuous.

The output equation is also a reduced-form gap relative to reference GDP. Its original
central path ends 3.45 percent above reference output. It is an overheating scenario,
not an estimated long-run loss of productive capacity or a household recession forecast.
No endogenous productivity, exchange-rate, banking, or capital-accumulation model supports
the richer economic events formerly described in the chronology.

\subsection*{Inflation persistence determines much of the apparent endgame}

The quarterly equations in the original central run are
\begin{align*}
x_{t+1} &= 0.80x_t-0.25\Delta r^{\rm real}_{t-1}
 -0.15s^{\rm private}_{t-1}+0.80d^P_t+d^I_t,\\
\pi_{t+1}-\pi^0_{t+1} &= 0.98(\pi_t-\pi^0_t)+0.06x_{t-1},\\
a_t &= 0.50a_{t-1}+0.50[1.50(\pi_t-\pi^0_t)+0.50x_t].
\end{align*}
Here $x$ is the output gap, $d^P$ the primary-deficit impulse from automatic stabilizers,
$d^I$ the spendable incremental-interest impulse, and $a$ the nominal policy-rate
adjustment. Rates and shares are decimals. Real tightening subtracts current inflation
rather than a separately modeled expectation. Positive output gaps reduce the primary
deficit through an automatic-stabilizer coefficient of 0.45.

Absent other feedback, persistence of 0.98 gives an inflation deviation a half-life of
about 8.6 years; 0.90 gives about 1.6 years. At a fixed output gap, the implied stationary
inflation deviation is $0.06x/(1-\rho_\pi)$: $3x$ at 0.98 and $0.6x$ at 0.90.
The full dynamic system is more complicated, but this comparison exposes a powerful
unestimated assumption. A four-percent reporting line is neither an equilibrium-stability
test nor a definition of fiscal dominance.

\section{New challenges to the central inflation result}

The review reruns the sequential engine with visible changes to the assumptions. These
alternatives are not newly estimated coefficients and carry no probability weights. Except
where indicated, they retain the 500-basis-point, forty-quarter Treasury stress and the
original interest allocation. All rows in Table~\ref{tab:review} use automatic financing
and omit the optional debt--yield feedback, allowing the macro channels to be compared.

\begin{table}[htbp]
\centering
\small
\setlength{\tabcolsep}{4pt}
\begin{tabularx}{\textwidth}{Yrrr}
\toprule
Change from original stress & Peak inflation & 2056 inflation & 2056 debt/GDP\\
\midrule
None: original coefficients & 8.04\% & 8.04\% & 255.61\%\\
Treasury premium 100 bp; private spread unchanged & 2.66\% & 2.66\% & 185.04\%\\
Inflation persistence 0.90 instead of 0.98 & 3.15\% & 2.89\% & 259.92\%\\
Policy-rate output coefficient 0.50 instead of 0.25 & 5.08\% & 5.08\% & 258.07\%\\
Both recipient spending propensities halved & 3.57\% & 3.57\% & 255.80\%\\
Half the Treasury premium also enters private credit & 7.01\% & 7.01\% & 250.90\%\\
Persistence 0.90, output coefficient 0.50, and credit pass-through &
2.82\% & 2.59\% & 256.84\%\\
\bottomrule
\end{tabularx}
\caption{Assumption sensitivity, 2026Q4--2056Q3. Credit pass-through adds 250 basis points
for forty quarters to the original eight-quarter, 200-basis-point private spread. It does
not model a full private yield curve or bond wealth losses.}
\label{tab:review}
\end{table}

\textbf{The principal finding is that large debt accumulation survives changes that remove
the high-inflation outcome.} The 0.90-persistence row never reaches four-percent inflation,
without reducing the original bondholder spending share. This directly limits the earlier
claim that spending above roughly ten percent necessarily defeats monetary control. That
threshold applied to a policy-rule grid holding the rest of the macro structure fixed.

The original high-spending scenarios also leave the range in which a linear gap model is
credible: one produces 41-percent inflation and another reaches the numerical 99-percent
rate bound. A complete cash ledger does not validate the macro interpretation of those
numbers. They identify extreme sensitivity and model boundaries.

\section{Why the 2039 financing date is not a discovered fiscal limit}

The earlier capacity experiment lets a common additional yield spread increase absorption
until a fixed maximum is reached. If $Q_t$ is gross issuance, $Q_t^0$ reference issuance
scaled to scenario GDP, and $\Delta y_t$ the nominal yield gap, the central condition is
\[
\frac{Q_t}{Q_t^0}\le
\min\left(1.5,\max\left[0,1+0.05\frac{\Delta y_t}{100\text{ bp}}\right]\right).
\]
At a required multiple above 1.5, the code stops regardless of the yield offered. With the
original macro coefficients and debt--yield feedback, this occurs in 2039Q3. Tightening
the maximum to 1.25 moves the boundary to 2034Q3; widening it to 2.0 avoids a boundary
through 2056. These dates are valid outputs of the supplied curves.

The curves do not establish a U.S. economic limit. Gross issuance includes repeated
refinancing of the same short-term principal, and reference financing includes Fed-held
debt as well as private holdings. A maximum multiple of that flow is not derived from
household wealth, intermediary balance sheets, collateral needs, net absorption, expected
returns, or the government's capacity to generate future surpluses. Changing issuance
maturities can change gross turnover without a corresponding change in long-run solvency.

Refinancing volume matters for liquidity and dealer intermediation, so it should remain
a diagnostic. A meaningful capacity model would distinguish reinvestment of principal,
net portfolio accumulation, maturities and duration risk, intermediary constraints, and
monetary support. The existing ceiling is a sensitivity parameter, not an empirically
identified barrier. CBO's fiscal-crisis analysis likewise supplies no universal debt ratio
or date at which investor confidence must fail \cite{cboCrisis}.

\section{What the eventual adjustments mean}

The policy conflict can resolve gradually or abruptly, and several adjustments can coexist.
The table describes economic mechanisms; it does not rank their probability.

\begin{center}
\small
\renewcommand{\arraystretch}{1.13}
\begin{tabularx}{\textwidth}{p{0.22\textwidth}YY}
\toprule
Adjustment & What changes & Who bears the cost or risk\\
\midrule
Fiscal reform or benefit limits & Primary borrowing falls; future budget backing improves. &
Taxpayers or recipients; incidence depends on the actual tax and spending provisions.\\
\addlinespace
Favorable growth or voluntary low returns & The tax base and desired liabilities expand
fast enough to accommodate borrowing. & No necessary crisis loss; the needed productivity,
saving, and safe-asset demand are uncertain.\\
\addlinespace
Inflation with accommodation & Nominal claims and insufficiently indexed payments buy
less; new financing must be consistent with the monetary regime. & Holders of exposed
nominal claims and households whose incomes lag prices.\\
\addlinespace
Financial repression & Regulation, subsidized absorption, or rate suppression lowers
financing costs relative to market alternatives. & Savers and institutions accepting
lower returns; possible costs to credit allocation.\\
\addlinespace
Payment delay or restructuring & The amount or timing of promised payments changes. &
Affected claimants, with potential losses transmitted through collateral and financial markets.\\
\bottomrule
\end{tabularx}
\end{center}

\subsection*{The Federal Reserve can alter financing conditions, not abolish the constraint}

The Fed can create dollar reserve balances and purchase securities in secondary markets.
Current arrangements exclude competitive bids by the Fed for new Treasury issuance;
SOMA rollovers are noncompetitive auction add-ons priced at the competitive stop-out yield
\cite{fedTreasuryPurchases,nyFedRollovers}. Dollar denomination gives the consolidated
public sector monetary options that a household or foreign-currency borrower lacks.
It does not give Treasury unlimited independent legal spending authority.

When a Fed purchase replaces a privately held Treasury with interest-bearing reserves,
the consolidated public sector has changed the form and maturity of its liability.
The Treasury security held inside the public sector must not be counted a second time
as an external claim. Reserve interest, remittances, and losses matter for financing cost.
Paying market rates on reserves does not make the debt free; moving toward liabilities
paying below-market returns relies on demand for money or on restrictions supporting them.
Real money demand is not an unlimited inflation-tax base.

Fiscal-dominance theory shows how maintaining a fiscal path without adequate backing can
limit monetary control, depending on how the authorities resolve the conflict. It does
not say that every rate increase is inflationary under every fiscal regime
\cite{andolfatto}. This simulator's interest-spending rule is not a structural model of
fiscal price-level determination, money demand, or expectations about a regime switch.

\subsection*{Inflation cannot be counted as permanent relief while its financing effects are ignored}

Unexpected inflation can reduce the real value of outstanding fixed nominal debt. TIPS
are indexed; bills and FRNs reprice quickly; future nominal yields can incorporate expected
inflation and risk. Indexed spending and tax responses also change the primary balance.
A one-time erosion of old debt does not repair a persistent flow imbalance by itself.
If nominal rates are held down while inflation rises, that is an additional monetary or
regulatory assumption that must be stated.

Thus a debt/GDP decline during high inflation is not automatically improved household
welfare or restored fiscal capacity. Nor does a higher price level imply that everyone's
real income falls by the same percentage. A household account requires wages, taxes,
benefits, asset holdings, liabilities, and indexation rules; this model does not contain
those household balance sheets.

\section{What remains useful in the conditional recovery}

The existing recovery run begins with the Treasury cohorts inherited from the central
2039Q2 capacity scenario. Debt is 178.82 percent of GDP, inflation 4.89 percent, and
annualized modeled interest 20.49 percent of GDP. These are highly stressed conditional
states, not a forecast of the economy in 2039.

The continuation assumes inflation peaks at 5.2 percent, an initial period of
$-1$-percent real growth, and a gradual decline in new bill and long-bond rates to
4.3 and 5.1 percent. Under these supplied paths, a permanent 5.48-percent-of-GDP primary
improvement, phased in over sixteen quarters, returns debt/GDP to 178.82 percent by 2056
and leaves it falling. Debt first peaks at 209.27 percent. The same favorable macro paths
without the fiscal improvement end at 265.14 percent.

This is useful accounting of the inherited coupon burden. The simulation does not show
that enacting the package would cause the assumed recession, inflation decline, or yield
normalization. Temporary Fed support, market confidence, and congressional agreement are
outside that recovery calculation. A seventeen-year endpoint is not an infinite-horizon
solvency proof. The 5.48-percent package is also not a directly comparable estimate of the
cost of waiting relative to the new 2.68-percent preventive calculation: the inherited
stocks, macro paths, horizons, and terminal targets differ.

\section{Research verdict and the work a predictive model would need}

The Treasury engine answers how existing securities reprice, how much financing a
specified path requires, and what fiscal adjustment meets a stated debt condition.
The new decomposition makes the no-shock imbalance explicit. Those are defensible
quantitative results. The confidence shock, recipient spending, inflation persistence,
absorption ceiling, and recovery environment have different evidentiary status.

A predictive extension should start with separate scheduled-benefits and payable-benefits
fiscal paths, then model how primary balances respond to debt and economic conditions.
It should measure aggregate interest transmission including private borrower exposure,
asset revaluation, taxation, and the consolidated Fed balance sheet. Investor demand
should respond to expected returns and fiscal backing, with liquidity constraints kept
distinct from long-run solvency. Expectations about policy changes, uncertainty in growth
and rates, and validation against historical monetary and fiscal responses are necessary
before assigning crisis probabilities. Adding one more unestimated crisis coefficient
would not resolve these identification problems.

The conclusion of this review is therefore sharper than a dated catastrophe story:
\textbf{continued deficits are possible; continued fiscal promises that outrun their
economic support are not indefinitely compatible with unchanged real payment terms.}
Under this project's baseline, the debt burden is still increasing substantially in 2056.
If favorable growth and demand for liabilities do not close the gap, an adjustment must
fall on taxes, spending, real creditor returns, or payments. It can occur before any
auction failure, and monetary accommodation is one conditional resolution. The evidence
does not identify an inevitable 2039 crisis, an unavoidable inflation spiral, or a single
final U.S. outcome.

\section*{Reproducibility and scope of this revision}

The review adds \path{src/debt_sim/sustainability.py} and a separate experiment. It leaves
the previous engine coefficients and historical scenario outputs intact. Run from the
repository root:
\begin{verbatim}
uv run python scripts/run_sustainability_review.py
uv run pytest
cd docs
latexmk -pdf -interaction=nonstopmode -halt-on-error \
  confidence_shock_review.tex
\end{verbatim}
The new paths, fiscal solutions, constant-policy examples, and exact debt decomposition
are \path{data/processed/sustainability_review_*.csv}. Scenario parameters are recorded in
\path{config/scenarios/sustainability_review.json}; Figure~\ref{fig:no-shock} reads
\path{docs/sustainability_review_plot_data.csv}. Tables report rounded values from those
outputs. The tests check a sustainable permanent-deficit example, the equal-rate/growth
case, and reconciliation of the cohort ledger including indexation and other financing.

The older conditional stress and recovery artifacts can be regenerated with
\path{scripts/run_debt_yield_feedback_experiment.py} and
\path{scripts/run_recovery_settlement_experiment.py}. The following appendix retains their
plots for comparison. It uses the original coefficients, including inflation persistence
0.98. No new macro calibration or probability model is claimed in this revision.

\clearpage
\appendix
\section{The earlier stress and recovery, retained as conditional illustrations}

The first two figures use the original central debt--yield scenario with automatic
financing through 2056. They do \emph{not} use the finite-capacity run that stops in 2039.
The third figure is a separate recovery calculation from that capacity run's last
successful quarter. These are different experiments, not one continuous forecast.
The recovery's macro and rate paths are imposed; its debt and interest paths are computed.

\begin{figure}[htbp]
\centering
\begin{tikzpicture}
\begin{groupplot}[
  group style={group size=2 by 1, horizontal sep=1.15cm},
  width=0.47\textwidth,
  height=0.31\textwidth,
  xmin=2026.5, xmax=2057,
  xtick={2027,2036,2046,2056},
  grid=major,
  tick label style={font=\small},
  label style={font=\small},
  title style={font=\small\bfseries},
  grid style={draw=gray!25},
  legend style={font=\scriptsize, draw=none, fill=white, fill opacity=0.85,
                text opacity=1}
]
\nextgroupplot[
  title={Pressure added to market rates},
  ylabel={Basis points},
  ymin=0, ymax=550,
  legend pos=south east
]
\addplot[policygold, very thick, const plot]
  table[col sep=comma, x=plot_year, y=treasury_premium_bp]
  {confidence_shock_central_plot_data.csv};
\addlegendentry{Treasury shock: input}
\addplot[quietgray, dashed, very thick, const plot, mark=*, mark size=1pt,
         mark indices={0,1,2,3,4,5,6,7}]
  table[col sep=comma, x=plot_year, y=private_credit_spread_bp]
  {confidence_shock_central_plot_data.csv};
\addlegendentry{Private-credit shock: input}
\addplot[stressred, very thick]
  table[col sep=comma, x=plot_year, y=debt_yield_long_bp]
  {confidence_shock_central_plot_data.csv};
\addlegendentry{Long-debt premium: result}

\nextgroupplot[
  title={Rates on newly issued debt},
  ylabel={Percent},
  ymin=0, ymax=22,
  legend pos=north west
]
\addplot[black, dashed, thick]
  table[col sep=comma, x=plot_year, y=policy_rate_pct]
  {confidence_shock_central_plot_data.csv};
\addlegendentry{Fed short-rate setting}
\addplot[policygold, very thick]
  table[col sep=comma, x=plot_year, y=bill_rate_pct]
  {confidence_shock_central_plot_data.csv};
\addlegendentry{Bills}
\addplot[stressred, very thick]
  table[col sep=comma, x=plot_year, y=long_rate_pct]
  {confidence_shock_central_plot_data.csv};
\addlegendentry{Long bonds}
\addplot[quietgray, densely dotted, thick, forget plot]
  coordinates {(2036.625,0) (2036.625,22)};
\end{groupplot}
\end{tikzpicture}
\caption{The rate sequence in the unresolved stress path.  The two rectangular paths in the left
panel are imposed stresses.  The rising red path and every rate in the right panel are
calculated before any recovery settlement.  The vertical dotted line marks the removal of the
five-point Treasury premium after 2036Q3.}
\label{fig:visual-rates}
\end{figure}
\FloatBarrier

\begin{figure}[htbp]
\centering
\begin{tikzpicture}
\begin{groupplot}[
  group style={group size=3 by 1, horizontal sep=0.82cm},
  width=0.315\textwidth,
  height=0.29\textwidth,
  xmin=2026.5, xmax=2057,
  xtick={2027,2041,2056},
  grid=major,
  tick label style={font=\scriptsize},
  label style={font=\small},
  title style={font=\scriptsize\bfseries},
  grid style={draw=gray!25},
  legend style={font=\scriptsize, draw=none, fill=white, fill opacity=0.85,
                text opacity=1}
]
\nextgroupplot[
  title={Interest entering demand},
  ylabel={Trillions per quarter},
  ymin=0, ymax=4.3
]
\addplot[policygold, very thick, fill=policygold, fill opacity=0.16]
  table[col sep=comma, x=plot_year, y=domestic_demand_trillions]
  {confidence_shock_central_plot_data.csv} \closedcycle;

\nextgroupplot[
  title={Output gap},
  ylabel={Percent},
  ymin=-0.3, ymax=7.3
]
\addplot[debtblue, very thick]
  table[col sep=comma, x=plot_year, y=output_gap_pct]
  {confidence_shock_central_plot_data.csv};
\addplot[quietgray, densely dotted, thick, forget plot]
  coordinates {(2026.5,0) (2057,0)};

\nextgroupplot[
  title={Inflation},
  ylabel={Annual percent},
  ymin=0, ymax=12.3,
  legend pos=north west
]
\addplot[quietgray, dashed, thick]
  table[col sep=comma, x=plot_year, y=reference_inflation_pct]
  {confidence_shock_central_plot_data.csv};
\addlegendentry{Reference}
\addplot[stressred, very thick]
  table[col sep=comma, x=plot_year, y=inflation_pct]
  {confidence_shock_central_plot_data.csv};
\addlegendentry{Stress path}
\addplot[black, densely dotted, thick, forget plot]
  coordinates {(2026.5,4) (2057,4)};
\end{groupplot}
\end{tikzpicture}
\caption{The fiscal--monetary feedback in pictures.  Higher rates create additional cash
interest; the assumed spendable share supports demand; output and inflation move above their
reference paths; and the Fed responds with rates that make the next refinancing round more
expensive.  The dotted line in the last panel is the paper's 4-percent early-warning level.}
\label{fig:visual-feedback}
\end{figure}
\FloatBarrier

\begin{figure}[htbp]
\centering
\begin{tikzpicture}
\begin{groupplot}[
  group style={group size=3 by 1, horizontal sep=0.83cm},
  width=0.315\textwidth,
  height=0.30\textwidth,
  xmin=2039.4, xmax=2057,
  xtick={2040,2044,2050,2056},
  grid=major,
  tick label style={font=\scriptsize},
  label style={font=\small},
  title style={font=\scriptsize\bfseries},
  grid style={draw=gray!25},
  legend style={font=\scriptsize, draw=none, fill=white, fill opacity=0.85,
                text opacity=1}
]
\nextgroupplot[
  title={Debt after the settlement},
  ylabel={Percent of GDP},
  ymin=165, ymax=275,
  legend pos=north west
]
\addplot[debtblue, very thick]
  table[col sep=comma, x=plot_year, y=debt_gdp_pct]
  {recovery_settlement_plot_data.csv};
\addlegendentry{Fiscal settlement}
\addplot[quietgray, dashed, thick]
  table[col sep=comma, x=plot_year, y=no_fiscal_debt_gdp_pct]
  {recovery_settlement_plot_data.csv};
\addlegendentry{Rates normalize only}

\nextgroupplot[
  title={The budget turns gradually},
  ylabel={Percent of GDP},
  ymin=0, ymax=23,
  legend pos=north east
]
\addplot[stressred, very thick]
  table[col sep=comma, x=plot_year, y=interest_gdp_pct]
  {recovery_settlement_plot_data.csv};
\addlegendentry{Interest expense}
\addplot[policygold, very thick]
  table[col sep=comma, x=plot_year, y=primary_balance_gdp_pct]
  {recovery_settlement_plot_data.csv};
\addlegendentry{Primary surplus}

\nextgroupplot[
  title={Declared normalization path},
  ylabel={Annual percent},
  ymin=0, ymax=10,
  legend pos=north east
]
\addplot[stressred, very thick]
  table[col sep=comma, x=plot_year, y=inflation_pct]
  {recovery_settlement_plot_data.csv};
\addlegendentry{Inflation}
\addplot[policygold, very thick]
  table[col sep=comma, x=plot_year, y=short_rate_pct]
  {recovery_settlement_plot_data.csv};
\addlegendentry{Bills}
\addplot[debtblue, dashed, thick]
  table[col sep=comma, x=plot_year, y=long_rate_pct]
  {recovery_settlement_plot_data.csv};
\addlegendentry{Long bonds}
\end{groupplot}
\end{tikzpicture}
\caption{The conditional recovery after the 2039 financing boundary.  The left panel
compares the solved fiscal settlement with the same macro and rate assumptions absent the
primary-balance improvement.  Inflation and marginal rates in the right panel are scenario
assumptions; the cohort engine calculates the two debt paths and the interest path.}
\label{fig:visual-recovery}
\end{figure}
\FloatBarrier

\begin{thebibliography}{99}

\bibitem{cbo2026}
Congressional Budget Office, \emph{The Budget and Economic Outlook: 2026 to 2036},
February 2026.  \url{https://www.cbo.gov/publication/62105}.

\bibitem{cboCrisis}
Congressional Budget Office, \emph{Federal Debt and the Risk of a Fiscal Crisis}, 2010.
\url{https://www.cbo.gov/publication/21625}.

\bibitem{tic}
U.S. Department of the Treasury, ``Table 5: Major Foreign Holders of Treasury Securities,''
June 2026.  \url{https://ticdata.treasury.gov/resource-center/data-chart-center/tic/Documents/slt_table5.html}.

\bibitem{fedDebtRates}
Abhik Bhatt, Anthony M. Diercks, Benjamin Eyal, and Arsenios Skaperdas, ``The Causal Effect
of Debt on Interest Rates,'' Finance and Economics Discussion Series 2026-031, Board of
Governors of the Federal Reserve System, May 2026.
\url{https://doi.org/10.17016/FEDS.2026.031}.

\bibitem{fedTreasuryPurchases}
Board of Governors of the Federal Reserve System, ``How Does the Federal Reserve's Buying
and Selling of Securities Relate to the Borrowing Decisions of the Federal Government?,''
updated July 2024.
\url{https://www.federalreserve.gov/faqs/how-does-the-federal-reserve-buying-and-selling-of-securities-relate-to-the-borrowing-decisions-of-the-federal-government.htm}.

\bibitem{nyFedRollovers}
Federal Reserve Bank of New York, ``FAQs: Treasury Rollovers,'' May 2020.
\url{https://www.newyorkfed.org/markets/treasury-rollover-faq-05062020}.

\bibitem{auclert}
Adrien Auclert, ``Monetary Policy and the Redistribution Channel,'' \emph{American Economic
Review}, vol. 109, no. 6, 2019, pp. 2333--2367.
\url{https://doi.org/10.1257/aer.20160137}.

\bibitem{bisDebt2026}
Bank for International Settlements, ``High Public Debt and Shifting Financial Markets:
Challenges for Central Banks,'' \emph{Annual Economic Report 2026}, chapter 2.
\url{https://www.bis.org/publ/arpdf/ar2026e2.htm}.

\bibitem{blanchard}
Olivier Blanchard, ``Public Debt and Low Interest Rates,'' \emph{American Economic Review},
109(4), 2019, pp. 1197--1229.
\url{https://www.aeaweb.org/articles?id=10.1257/aer.109.4.1197}.

\bibitem{andolfatto}
David Andolfatto, ``Is It Time for Some Unpleasant Monetarist Arithmetic?,''
\emph{Federal Reserve Bank of St. Louis Review}, 103(3), 2021, pp. 315--332.
\url{https://www.stlouisfed.org/publications/review/2021/05/26/is-it-time-for-some-unpleasant-monetarist-arithmetic}.

\bibitem{bernankeSustainability}
Ben S. Bernanke, ``Achieving Fiscal Sustainability,'' National Commission on Fiscal
Responsibility and Reform, April 27, 2010.
\url{https://www.federalreserve.gov/newsevents/speech/bernanke20100427a.htm}.

\end{thebibliography}

\end{document}
